%% Direct contextual first-order soundness, retaining possibly empty sorts.

\section{Contextual first-order interpretation and soundness}

\begin{annexitem}{annex:p10.quant-fol.1}{thm:p10-doctrine-substitution}{Substitution as composition and reindexing}
\begin{proof}
All terms, formulas, and contexts in one substitution instance are finite.
Fresh bound names of the same sorts can therefore be chosen before
substitution, outside all names in the substituting terms and retained
contexts. Structural induction proves well-formedness: variables are
replaced by terms of their declared sorts, function symbols retain
their typed argument lists, and formula constructors retain their
required arguments. At a quantifier, the fresh variable is retained
unchanged while the other variables are substituted.

For term interpretations, the variable case is the projection equation
of the tuple map \(s\). The function-symbol case uses induction on
its arguments and
\(\langle f_i\rangle\circ s=\langle f_i\circ s\rangle\), including
the unique leading terminal component.
The nullary case is \(!\circ s=!\).
Thus \(\llbracket M[\mathbf s]\rrbracket=\llbracket M\rrbracket\circ s\).

For atomic relations, the substituted argument tuple is the old tuple
composed with \(s\). Functoriality of reindexing proves the formula
identity. This includes nullary relations, reindexed along the unique
map to \(1\).
For equality the pair of substituted term arrows is
\(\langle\llbracket M\rrbracket,\llbracket N\rrbracket\rangle\circ s\);
reindexing the fixed predicate \(\delta_{A_\sigma}\) gives the same result.
The false, conjunction, disjunction, and implication cases use
preservation of those Heyting operations by \(s^*\).

For a quantified formula take \(I=I_\Gamma\), \(J=I_\Delta\),
\(A=A_\sigma\), with its fresh bound variable appended last.
The lifted substitution is interpreted by
\(s\times1_A:J\times A\to I\times A\): its old components are
weakened by the first projection, by the term substitution case already
proved, and its new component is the last
variable projection. For \(b=\llbracket\varphi(x)\rrbracket_{\Gamma,x:\sigma}\),
the induction hypothesis and the two Beck--Chevalley equations give
\[
\begin{aligned}
 \llbracket(\exists_\sigma(\varphi))[\mathbf s]\rrbracket_\Delta
   &=\exists_{\pi_{J,A}}(s\times1_A)^*b
     =s^*\exists_{\pi_{I,A}}b,\\
 \llbracket(\forall_\sigma(\varphi))[\mathbf s]\rrbracket_\Delta
   &=\forall_{\pi_{J,A}}(s\times1_A)^*b
     =s^*\forall_{\pi_{I,A}}b.
\end{aligned}
\]
Thus each quantifier has the required substitution law.
Consistent renaming of a bound variable retains the same sort and
final projection, so the same induction proves bound alpha-invariance.
Freshening choices consequently give the same interpretation.

Variable substitutions that select old declarations give weakening
and context permutations. The associated arrows are the corresponding
product projections and permutation isomorphisms. Hence weakening a
formula independent of a final \(x:\sigma\) gives
\(\llbracket\varphi\rrbracket_{\Gamma,x:\sigma}
 =\pi_{I_\Gamma,A_\sigma}^*\llbracket\varphi\rrbracket_\Gamma\).
No assertion that this reindexing reflects order is needed.
\end{proof}
\end{annexitem}

\begin{annexitem}{annex:p10.quant-fol.2}{thm:p10-doctrine-soundness}{Soundness of every displayed contextual rule}
\begin{proof}
Fix an interpretation satisfying the theory axioms. We check every rule.

\paragraph{Equational rules and theory axioms.}
The two axiom rules are valid by assumption.
Eq-Refl, Eq-Sym, and Eq-Trans are the corresponding laws of equality
of arrows. In Eq-Subst the equation \(M=M'\) in the prefix context
makes the two tuple maps inserting \(M\) or \(M'\) into
\(\Gamma,x:\sigma,\Delta\) equal. Call that common map \(s\).
The other equation has equal term arrows \(n=n'\).
Term substitution gives \(n\circ s=n'\circ s\), as required.
This also validates operation congruence and weakening whenever
derived from these equational rules.

\paragraph{Structural sequent rules.}
In any fixed context write \(a=\bigwedge\llbracket\Phi\rrbracket\),
\(a'=\bigwedge\llbracket\Phi'\rrbracket\), and use lower-case letters
for conclusion and displayed formula values.
Id is \(a\wedge b\leq b\).
For Cut, \(a\leq b\) and \(a'\wedge b\leq c\) imply
\(a\wedge a'\leq a'\wedge b\leq c\).
Contraction and exchange preserve the antecedent meet by idempotence
and commutativity.
Context-Weaken reindexes the inequality along the map forgetting the
new declarations; reindexing preserves the finite antecedent meet.
For Subst, the equational premise again makes the inserting maps
\(s_M,s_{M'}\) equal. Reindex the old sequent along this common map.
Formula substitution identifies its antecedent with \(\Phi(M)\)
and its conclusion with \(\psi(M')\), proving the displayed rule.
Adding assumptions is also admissible: Cut a derivation of
\(\Phi\vdash\psi\) with Id for \(\Phi',\psi\vdash\psi\).

\paragraph{Propositional rules.}
False-Elim uses \(a\leq\bot\leq c\).
For conjunction introduction, \(a\leq b\) and \(a'\leq c\) give
\(a\wedge a'\leq b\wedge c\); each elimination uses a meet projection.
The two disjunction introductions use the join inclusions.
For disjunction elimination the premises are
\(a\wedge b\leq d\), \(a'\wedge c\leq d\), and
\(a''\leq b\vee c\).
Put \(r=a\wedge a'\wedge a''\).
Then \(r\leq b\vee c\), and distributivity gives
\(r=(r\wedge b)\vee(r\wedge c)\leq d\).
Implication introduction is the Heyting adjunction
\(a\wedge b\leq c\) iff \(a\leq b\Rightarrow c\).
For elimination, \(a\leq b\Rightarrow c\) and \(a'\leq b\)
give \(a\wedge a'\leq(b\Rightarrow c)\wedge b\leq c\).

\paragraph{Universal quantification.}
Let \(I=I_\Gamma\), \(A=A_\sigma\), and \(\pi:I\times A\to I\).
In universal introduction, the antecedent is independent of the
new variable. The premise is therefore \(\pi^*a\leq b\), and
the right adjunction gives \(a\leq\forall_\pi b\).
For elimination, the term \(M\) gives the section
\(s=\langle1_I,\llbracket M\rrbracket\rangle:I\to I\times A\).
Reindex the counit \(\pi^*\forall_\pi b\leq b\) along \(s\).
Since \(\pi s=1_I\), this gives
\(\forall_\pi b\leq s^*b=\llbracket\varphi(M)\rrbracket_\Gamma\).
Composing with the premise inequality proves elimination.
Existence of this section is supplied by the rule's well-sorted term,
not by a general assumption that the sort is inhabited.

\paragraph{Existential quantification.}
Reindex the unit \(b\leq\pi^*\exists_\pi b\) along the same term
section to get \(s^*b\leq\exists_\pi b\); this proves introduction.
In elimination the second premise has the form
\(\pi^*a'\wedge b\leq\pi^*c\), because both its undischarged
assumptions and conclusion are independent of the eigenvariable.
The left adjunction and Theorem~\ref{thm:p10-doctrine-frobenius} give
\[
 a'\wedge\exists_\pi b
    =\exists_\pi(\pi^*a'\wedge b)\leq c.
\]
The first premise \(a\leq\exists_\pi b\) now gives
\(a\wedge a'\leq c\), exactly the concatenated hypothesis list.

\paragraph{Equality formulas.}
For equality introduction, Theorem~\ref{thm:p10-doctrine-equality-laws}
gives \(\Delta_A^*\delta_A=\top_A\).
Reindex along \(\llbracket M\rrbracket:I\to A\); the equality
\(\langle M,M\rangle=\Delta_A\circ M\) gives
\(\llbracket M=_\sigma M\rrbracket=\top_I\), above any antecedent.

For equality elimination let
\(b=\llbracket\psi(x,y)\rrbracket_{\Gamma,x:\sigma,y:\sigma}\)
and \(a'=\bigwedge\llbracket\Phi'\rrbracket_\Gamma\).
Use the left-associated maps \(e,p,r\) of
Definition~\ref{def:p10-doctrine-equality}.
By substitution, the second premise is \(\pi^*a'\leq e^*b\).
The parameterized equality adjunction gives
\[
                   p^*\pi^*a'\wedge r^*\delta_A\leq b.
\]
Reindex along
\(h=\langle\langle1_I,\llbracket M\rrbracket\rangle,
                         \llbracket N\rrbracket\rangle:I\to(I\times A)\times A\).
Since \(\pi p h=1_I\) and
\(r h=\langle\llbracket M\rrbracket,\llbracket N\rrbracket\rangle\),
this yields
\[
 a'\wedge\llbracket M=_\sigma N\rrbracket_\Gamma
                      \leq\llbracket\psi(M,N)\rrbracket_\Gamma.
\]
The first premise bounds its assumption meet by the equality
predicate, proving the required conclusion with both hypothesis lists.
No reflection from a predicate equality to equality of base arrows
has entered either equality-formula proof.

Induction on the finite derivation tree now proves soundness.
Every leaf that is a theory axiom is satisfied by hypothesis; the
finite tree contains only finitely many such leaves even if
\(\mathcal T\) is an arbitrary infinite set.
\end{proof}
\end{annexitem}

\begin{annexitem}{annex:p10.quant-fol.3}{thm:p10-doctrine-tarski}{Finite assignments and the inhabitance boundary}
\begin{proof}
In Set, the chosen leading-terminal product \(I_\Gamma\) has the
canonical bijection to the set of sorted assignments to the finitely
many variables of \(\Gamma\): read its variable projections.
Its inverse inserts the unique terminal coordinate and pairs the
assigned values in the declared order.
For the empty context both are singleton sets.
Transport each function and relation interpretation through the
corresponding argument-product bijection. This includes nullary
functions and nullary relations.

Induction on a term identifies its arrow value at an assignment with
ordinary recursive term evaluation. For an atomic relation, reindexing
is inverse image, so membership is exactly truth of that relation on
the evaluated arguments. For equality, the powerset equality predicate
is the diagonal, giving equality of the two term values.
False is the empty set, conjunction is intersection, disjunction is
union, and the Boolean Heyting implication is
\((I_\Gamma\setminus B)\cup C\). These are the classical Tarskian
connective clauses.

For a parameter assignment \(a\), the fibre of
\(I_\Gamma\times A_\sigma\to I_\Gamma\) is its possible values
for the new variable. Existential quantification asserts that some
such value satisfies the body; universal quantification asserts that
every such value does. These are exactly the powerset adjoints.
The induction therefore proves the stated agreement for all formulas.
Equality of term arrows is pointwise equality at all assignments,
and inclusion of the antecedent intersection in the conclusion set
is precisely the universal sequent condition.
Assignments involve finite contexts; no simultaneous choice of values
for all variables or all sorts of an infinite signature is required.

If every sort carrier is nonempty, these are the usual nonempty-sort
clauses. For an unused variable of such a sort, each assignment to
the shorter context has an extension: choose one value of that sort.
This is one existential choice for the given assignment, not a global
family of witnesses. Thus its projection is surjective, so inverse
image reflects inclusion and precomposition reflects equality of term
functions in Set. Equivalently, for the predicate assertion,
\(\exists_\pi\top=\top\), and
Theorem~\ref{thm:p10-doctrine-inhabited} proves the corresponding
order-reflection statement in any doctrine with inhabited sorts.
This abstract consequence concerns predicate sequents only; the
doctrine does not in general make the projection epic in its base.
Apply the appropriate statement one variable at a time to a finite
unused context suffix.

For an empty sort \(A_\sigma\), the product with a nonempty parameter
set has empty fibres. The existential value of truth is empty and
the universal value of false is the whole parameter set.
In particular the empty context is a singleton, so
\(\exists_\sigma((x)(x=_\sigma x))\) is false there.
The context containing \(x:\sigma\) is empty, where every sequent
inequality is satisfied, including the one with false conclusion.
Removing that declaration would turn it into a false inequality in
the singleton context. This proves the claimed boundary.

Finally, generic doctrine equality is not assumed to reflect arrows.
For example, take Set as base and the one-element Heyting algebra
as every fibre, with all reindexing and adjoints the unique maps.
The equality axioms hold, and every equality formula is true, even
for two distinct constant arrows into a two-element set.
The displayed calculus has no rule converting that fact into an
external term equation. The soundness theorem therefore applies to
this doctrine as well.
\end{proof}
\end{annexitem}
